\section{Roadmap to Unconditional Proof: Closing the Final Gaps}
\label{app:future-roadmap}

%=============================================================================
% THE FINAL FRONTIER
% This appendix outlines the specific mathematical programs required to 
% turn the "Hypotheses" of this paper into "Theorems".
%=============================================================================

\textbf{Note: This section has been revised in accordance with the Rigorization Addendum (Appendix \ref{app:rigor-addendum}).}

This document has reduced the Yang-Mills Mass Gap problem to two precise, isolated mathematical hypotheses. The "Roadmap to Unconditional Proof" is now defined by the proof of the following two hypotheses:

\begin{enumerate}
    \item \textbf{Hypothesis RG-Stability:} Uniform control of the effective action renormalization group flow.
    \item \textbf{Hypothesis Tightness:} Existence of the continuum limit in a suitable distribution space.
\end{enumerate}

This appendix provides the precise formulation of these hypotheses.

%=============================================================================
\subsection{Hypothesis 1: RG-Stability}
%=============================================================================

\begin{hypothesis}[RG-Stability]
\label{hyp:rg-stability}
Let $S_{eff}^{(k)}$ be the effective action at scale $k$. There exist constants $C_1, C_2, \kappa > 0$ such that for all $k$:
\begin{enumerate}
    \item (Convexity) $\nabla^2 S_{eff}^{(k)} \ge -C_1$ in the gauge-fixed subspace.
    \item (Decay) $|\frac{\partial^2 S_{eff}^{(k)}}{\partial \phi_x \partial \phi_y}| \le C_2 e^{-\kappa |x-y|}$.
    \item (Large Field) $S_{eff}^{(k)}(\phi) \ge c \|\phi\|^2 - D$.
\end{enumerate}
\end{hypothesis}

\textbf{The Task:} Prove that the renormalization group flow for 4D non-abelian Yang-Mills theory satisfies these bounds uniformly in the scale $k$. This requires a rigorous implementation of the Balaban/Gawedzki block-spin transformation with large-field stability.

%=============================================================================
\subsection{Hypothesis 2: Geometric Tightness}
%=============================================================================

\begin{hypothesis}[Tightness in $H^{-s}$]
\label{hyp:tightness-sobolev}
The sequence of lattice gauge fields $A^{(a)}$, when embedded in the continuum distribution space, satisfies:
\begin{equation}
\sup_{a \to 0} \mathbb{E} \left[ \| A^{(a)} \|_{H^{-s}}^p \right] < \infty
\end{equation}
for some $s > 2$ and $p \ge 1$. This ensures the existence of a subsequential limit in the space of distributions.
\end{hypothesis}

\textbf{The Task:} Prove this tightness bound. This replaces the heuristic $a^{2p}$ scaling assumption with a mathematically sound condition for distribution-valued fields.

%=============================================================================
\subsection{Conclusion}
%=============================================================================

The proof of these two hypotheses constitutes the remaining work for an unconditional Clay Millennium Prize solution.

\begin{table}[h]
\centering
\begin{tabular}{|l|l|l|}
\hline
\textbf{Hypothesis} & \textbf{Physical Meaning} & \textbf{Mathematical Challenge} \\
\hline
\textbf{RG Stability} & Effective interactions remain local & Controlling the non-abelian polymer expansion \\
\hline
\textbf{Geometric Tightness} & Fields don't become "too rough" & Renormalizing the 4D gauge SPDE (Landau pole) \\
\hline
\end{tabular}
\caption{Summary of Remaining Hard Problems}
\end{table}

Solving either of these would be a major breakthrough in its own right, likely worth a Fields Medal. The framework provided in this paper ensures that \textit{only} these specific problems need to be solved to claim the Clay Prize.
